\section{التركيب الفيزيائي لمحركات السيرفو}
\label{sec:physical_structure}

تعتمد محركات السيرفو الحديثة، وتحديداً المحركات التزامنية ذات المغناطيس الدائم (\textenglish{PMSM - Permanent Magnet Synchronous Motors})، على بنية هندسية دقيقة تجمع بين الكفاءة الكهرومغناطيسية وانخفاض عزم القصور الذاتي الميكانيكي (\textenglish{Low Inertia}). هذا التصميم المدمج هو ما يمنحها القدرة على التسارع والتباطؤ بشكل شبه لحظي.

يتكون محرك السيرفو من جزأين فيزيائيين رئيسيين:

\subsection{العضو الثابت (\textenglish{Stator})}
يمثل الهيكل الخارجي الثابت للمحرك، ويتألف من مجموعة من الرقائق المعدنية (\textenglish{Laminations}) المصنوعة من صلب السليكون، وذلك لتقليل المفاقيد المغناطيسية الناتجة عن التيارات الدوامية (\textenglish{Eddy Currents}). تُلف أسلاك نحاسية معزولة داخل مجاري العضو الثابت لتشكيل ملفات ثلاثية الطور (\textenglish{3-Phase Windings}). عند مرور تيار كهربائي قادم من المشغل (\textenglish{Servo Drive}) في هذه الملفات، يتولد مجال مغناطيسي دوار يعبر الفجوة الهوائية باتجاه العضو الدوار.

\subsection{العضو الدوار (\textenglish{Rotor})}
على عكس المحركات الحثية التقليدية التي تعتمد على تحريض التيار في العضو الدوار، يُصنع العضو الدوار في محركات السيرفو الحديثة من مغانط دائمة قوية جداً من مواد أرضية نادرة، مثل النيوديميوم (\textenglish{Neodymium}). يحقق هذا التصميم ميزتين جوهريتين:
\begin{itemize}
    \item \textbf{غياب المفاقيد النحاسية في العضو الدوار:} لعدم وجود ملفات نحاسية فيه، مما يقلل بشكل كبير من ارتفاع حرارة المحرك ويزيد من كفاءته.
    \item \textbf{انخفاض عزم القصور الذاتي:} نظراً لصغر حجم وخفة وزن المغانط الدائمة مقارنة بالملفات التقليدية، مما يسمح للمحرك بالاستجابة الديناميكية السريعة جداً لتغيرات السرعة والاتجاه المطلوبة من المتحكم.
\end{itemize}

\subsection{مبدأ توليد العزم الكهرومغناطيسي}
ينشأ العزم الميكانيكي (\textenglish{Torque}) نتيجة التجاذب والتنافر المستمر بين المجال المغناطيسي الدوار الناشئ في العضو الثابت، والمجال المغناطيسي الثابت للمغانط في العضو الدوار. يمكن التعبير عن العزم الميكانيكي المتولد رياضياً باستخدام المعادلة المبسطة التالية:
$$T = K_t \cdot I_q$$
حيث:
\begin{itemize}
    \item $T$: العزم الميكانيكي المتولد بوحدة نيوتن-متر (\textenglish{N.m}).
    \item $K_t$: ثابت العزم الخاص بالمحرك (\textenglish{Torque Constant})، وهو قيمة فيزيائية تحددها الشركة المصنعة بناءً على قوة المغانط وعدد لفات النحاس.
    \item $I_q$: التيار الفعال (العمودي) المار في ملفات العضو الثابت، والمتحكم فعلياً في توليد العزم.
\end{itemize}